\section{Related Work}

\paragraph{LLM-as-a-judge and multimodal reward evaluation.}
Large language models are increasingly used as automatic evaluators of generated responses. Early work established the LLM-as-a-judge paradigm and examined its agreement with human preferences, while also documenting position, verbosity, and self-enhancement biases \citep{zheng2023judging}. Subsequent benchmarks evaluate judges on challenging response pairs, test-time response selection, and critique-based refinement \citep{tan2025judgebench,zhou2025evaluating}. These studies show that judge performance depends on the task, candidate distribution, and evaluation protocol, motivating direct evaluation of the judges rather than assuming that a capable generator is also a reliable evaluator.

This paradigm has recently been extended to multimodal inputs. MLLM-as-a-Judge evaluates multimodal judges through scoring, pairwise comparison, and batch-ranking tasks, revealing substantial variation across evaluation formats \citep{chen2024mllm}. MJ-Bench studies multimodal judges for text-to-image generation across dimensions including alignment, safety, image quality, bias, and composition \citep{chen2025mjbench}, while Multimodal RewardBench 2 extends evaluation to image generation, image editing, interleaved generation, and multimodal reasoning \citep{hu2026multimodal}. Complementary work develops increasingly general or capable multimodal reward models, including Omni-Reward, R1-Reward, and BaseReward \citep{jin2026omnireward,zhang2026r1reward,zhang2026basereward}. These evaluations primarily characterize performance on fixed examples using accuracy, ranking agreement, or correlation with human preferences. RewardLens studies a different question: whether judges assigned similar factor-specific static preference scores also behave similarly when the visual evidence governing the correct preference is changed.

\paragraph{Behavioral and controlled evaluation.}
Aggregate held-out accuracy can conceal systematic behavioral failures. CheckList formalizes behavioral testing through targeted capabilities and input transformations, including invariance tests for meaning-preserving changes and directional tests for changes that should alter a prediction \citep{ribeiro2020beyond}. More broadly, shortcut learning describes how models can perform well on conventional benchmarks by exploiting predictive but non-causal features that fail under carefully chosen test conditions \citep{geirhos2020shortcut}. RewardLens adopts this behavioral perspective for multimodal judges by separately testing when a preference should change and when it should remain invariant.

Controlled contrast has a long history in vision-language evaluation. VQA v2 pairs the same question with semantically similar images that require different answers, reducing the usefulness of language-only priors and forcing predictions to depend on visual evidence \citep{goyal2017making}. VALSE uses validated foils to test grounding of existence, cardinality, attributes, actions, and spatial relations \citep{parcalabescu2022valse}. Winoground, CREPE, and SugarCrepe construct contrastive examples or hard negatives to diagnose compositional vision-language understanding \citep{thrush2022winoground,ma2023crepe,hsieh2023sugarcrepe}. More recent work uses visual programming to generate verifiable counterfactual multimodal examples \citep{gao2025counterfactual}. These benchmarks establish the value of controlled examples for isolating particular capabilities. However, they primarily evaluate answer models or vision-language representations, rather than pairwise judges, and do not ask how much behavioral variation remains after judges are matched on an independent static preference score.

\paragraph{Perturbation-based audits of multimodal judges.}
Recent studies directly examine whether multimodal reward models and judges rely on appropriate visual evidence. Multimodal reward models can exploit unimodal, particularly text-only, spurious correlations that support in-distribution accuracy but impair out-of-distribution generalization \citep{li2025devil}. Visual manipulations can also systematically inflate the scores assigned by LVLM judges in text-to-image evaluation \citep{hwang2025fooling}. MM-JudgeBias evaluates compositional biases through controlled perturbations to the query, image, and response, distinguishing sensitivity to relevant evidence from stability under irrelevant perturbations \citep{lee2026mmjudgebias}. Perceptual Judgment Bias similarly studies cases in which visual evidence conflicts with plausible textual cues and proposes perturbation-based training to improve perceptual fidelity \citep{park2026mitigating}. Multi-Crit examines another form of behavioral flexibility: whether multimodal judges can change their decisions appropriately when the requested evaluation criterion changes \citep{xiong2026multicrit}.

VisualFLIP is the closest prior work at the intervention level: it holds the question fixed, minimally changes task-critical visual evidence, and measures whether a multimodal prediction updates when the gold answer changes \citep{zhu2026visualflip}. RewardLens does not claim the first use of visual perturbations or prediction-update tests. It introduces a different comparison and measurement design. First, it measures factor-specific static preference accuracy on an independent natural-image set and studies the residual intervention-behavior differences among judges matched by that score. Second, each audit triplet includes both a relevant edit that reverses the gold preference and an irrelevant edit that preserves it, while holding the question, candidate responses, candidate order, camera, and rendering conditions fixed. Third, Relevant Adaptation and Irrelevant Invariance are conditioned on a correct base decision, separating selective updating from generic prediction instability. RewardLens therefore complements static judge benchmarks and prior perturbation studies; it characterizes observable intervention behavior rather than claiming to identify a judge's internal reasoning or visual-grounding mechanism.
